
\documentclass[11pt,a4paper]{article}
\usepackage[margin=1in]{geometry}
\usepackage{booktabs}
\usepackage{longtable}
\usepackage{hyperref}
\usepackage{url}
\Urlmuskip=0mu plus 1mu
\usepackage{xcolor}
\usepackage{amsmath}
\usepackage{enumitem}
\usepackage[T1]{fontenc}
\usepackage[utf8]{inputenc}
\title{Semantic Duplicate and Near-Plagiarism Detection Engine\\Technical Report}
\author{Data Mining Project 3}
\date{\today}
\begin{document}
\maketitle

\section{Overview}
This repository implements a command-line plagiarism and near-duplicate detection engine. It follows two algorithmic paths: (1) word shingling, MinHash, and locality sensitive hashing (LSH), and (2) TF-IDF weighted SimHash. The goal is a transparent, reproducible implementation rather than a wrapper around ready-made MinHash, LSH, or SimHash libraries.

\section{Preprocessing}
Each input document is normalized using Unicode NFKC, lowercasing, removal of unnecessary punctuation, whitespace compaction, English and Persian stopword filtering, and Persian/Arabic character unification. The normalized text is tokenized and converted to word shingles. The default shingle size is 3. Empty texts, very short texts, and unusual characters are handled without crashing: empty documents produce an empty shingle set, while very short documents become one short shingle.

\section{Exact Jaccard Similarity}
For two shingle sets $A$ and $B$, exact similarity is computed as:
\[
J(A,B)=\frac{|A\cap B|}{|A\cup B|}.
\]
Computing exact Jaccard for every pair in a collection of $n$ documents requires $n(n-1)/2$ comparisons, which is $O(n^2)$. This is the motivation for MinHash and LSH.

\section{MinHash}
MinHash signatures are implemented from scratch. Each shingle is mapped to a stable 64-bit hash. A deterministic family of universal hash functions is then applied to compute the minimum value for every signature row. The default signature length is 128. The approximate similarity of two documents is the fraction of equal positions in their signatures.

\section{LSH Candidate Generation}
The MinHash signature is split into bands. Each band is hashed into a bucket. Two documents become candidates when they share at least one bucket in at least one band. This avoids direct comparison of all pairs. On the included sample corpus, LSH reduces comparisons from 21 all-pairs comparisons to 3 candidate comparisons, a reduction ratio of approximately 0.857.

\section{TF-IDF Weighted SimHash}
For the second path, each document is represented by token counts. Smoothed IDF is computed over the corpus or pair being evaluated. For every token, the TF-IDF weight is added or subtracted from a 64-dimensional vector according to the 64-bit token hash. The final bit is set to one if the corresponding vector dimension is non-negative. Similar documents should have low Hamming distance and therefore high similarity $1 - d_H/64$.

\section{CLI Commands}
The project exposes the three required commands plus one PAN-PC-11 preparation helper:
\begin{enumerate}[label=\arabic*.]
  \item \texttt{compare}: compare two text files and write JSON scores.
  \item \texttt{corpus}: search a folder of documents using MinHash+LSH candidates and write candidate pairs.
  \item \texttt{pairs}: evaluate labelled document pairs and write metrics.
  \item \texttt{prepare-pan11}: convert PAN-PC-11 XML annotations into a labelled pair CSV for the \texttt{pairs} command.
\end{enumerate}

\section{Dataset}
The project brief recommends the PAN Plagiarism Corpus 2011 (PAN-PC-11), DOI \path{10.5281/zenodo.3250095}. The corpus is distributed on Zenodo as a two-part RAR archive. Because the raw corpus is 1.7 GB, it is intentionally not committed to the repository.

\begin{center}
\scriptsize
\begin{tabular}{p{0.40\linewidth}p{0.10\linewidth}p{0.34\linewidth}}
\toprule
File & Size & MD5 \\
\midrule
\path{pan-plagiarism-corpus-2011.part1.rar} & 1.0 GB & \texttt{b2930f859497dd48ba5bb606d3f4a4f3} \\
\path{pan-plagiarism-corpus-2011.part2.rar} & 703.9 MB & \texttt{b23d86c17a47d2bfbdc4c314ea5810df} \\
\bottomrule
\end{tabular}
\end{center}

A downloader helper is included at \path{scripts/download_pan11.sh}, and dataset notes are included at \path{docs/PAN11_DATASET.md}. The repository also includes a small reproducible educational corpus under \path{data/sample_corpus/} and a small labelled pair CSV under \path{data/raw/train.csv} so the CLI, tests, and report outputs can be reproduced without downloading the full PAN archive.

For PAN-PC-11, the \texttt{prepare-pan11} command parses XML annotations and builds \path{data/processed/pan11_pairs.csv}. Positive examples are created from annotated suspicious/source plagiarism passages. Negative examples are sampled from unrelated source documents with a fixed random seed. This keeps the raw PAN archive outside GitHub while producing a reproducible pairwise file for the required \texttt{pairs} evaluator.

\section{Experimental Results}
The following results are generated by the included labelled sample pairs using default parameters: shingle size 3, 128 MinHash permutations, 32 LSH bands, MinHash threshold 0.25, and SimHash threshold 0.75.

\begin{center}
\begin{tabular}{lrrrrrrrr}
\toprule
Method & Acc. & Prec. & Recall & F1 & TP & FP & TN & FN \\
\midrule
MinHash\_similarity\_threshold & 0.900 & 1.000 & 0.800 & 0.889 & 4 & 0 & 5 & 1 \\
TFIDF\_weighted\_SimHash & 0.800 & 1.000 & 0.600 & 0.750 & 3 & 0 & 5 & 2 \\
\bottomrule
\end{tabular}
\end{center}

\section{Error Analysis}
The sample data intentionally includes both almost copied texts and paraphrased texts. False negatives mainly occur when a pair is semantically similar but does not preserve enough contiguous word shingles, or when different important terms lead to different SimHash bits. Examples from the current run:
\begin{itemize}

  \item Pair 3: label=1, MinHash=0.328, SimHash=0.594. The text is semantically related but has limited exact shingle overlap or different key tokens.
  \item Pair 4: label=1, MinHash=0.078, SimHash=0.719. The text is semantically related but has limited exact shingle overlap or different key tokens.
\end{itemize}

\section{Reproducibility}
Install and run:
\begin{verbatim}
python -m venv .venv
source .venv/bin/activate
pip install -r requirements.txt -e .
python -m plagiarism_engine.cli compare \
  --file-a data/sample_corpus/doc_01.txt \
  --file-b data/sample_corpus/doc_02.txt \
  --output outputs/two_file_compare.json
python -m plagiarism_engine.cli corpus \
  --data data/sample_corpus \
  --threshold 0.25 \
  --shingle-size 3 \
  --output outputs/candidates.csv
python -m plagiarism_engine.cli pairs \
  --pairs data/raw/train.csv \
  --text-col-a question1 \
  --text-col-b question2 \
  --label-col is_duplicate \
  --limit 5000 \
  --output outputs/metrics.csv
python -m plagiarism_engine.cli prepare-pan11 \
  --pan-root data/raw/pan11 \
  --output data/processed/pan11_pairs.csv \
  --max-positive 5000 \
  --negatives-per-positive 1 \
  --segment-mode annotated
python -m plagiarism_engine.cli pairs \
  --pairs data/processed/pan11_pairs.csv \
  --text-col-a text_a \
  --text-col-b text_b \
  --label-col label \
  --limit 5000 \
  --output outputs/metrics_pan11.csv
pytest -q
\end{verbatim}

\section{Conclusion}
The implementation satisfies the project requirements with modular source files, a command-line interface, explicit preprocessing, from-scratch MinHash and SimHash, LSH candidate reduction, reproducible outputs, unit tests, a dataset downloader note, PAN-PC-11 XML conversion, and a technical report. The bundled executable data is intentionally small for GitHub compatibility, while the PAN workflow can be run locally after downloading and extracting the official Zenodo archive.
\end{document}
